\documentclass{article}
\usepackage[utf8]{inputenc}
\usepackage[T1]{fontenc}
\usepackage{geometry}
\usepackage{graphicx}
\usepackage{float}
\usepackage{hyperref}
\usepackage[english,indonesian]{babel}

\title{Optimasi Pemrosesan Template LaTeX Menggunakan Pendekatan Berbasis AI}
\author{Academic Support Agent}
\date{\today}

\begin{document}
\maketitle

\begin{abstract}

\end{abstract}

\chapter{PENDAHULUAN}


\chapter{TINJAUAN PUSTAKA}


% Original content (commented out):
% \documentclass{article}
\usepackage{commands}

\title{
	\Huge{\bf{SUBJECT - TITLE}}
}
\author{
NAME \\
	\texttt{EMAIL} \\
	STUDENT NUMBER
}
\date{
	MONTH DAY\ts{th}, YEAR
	\vspace{1cm}
	\hrule
	\vspace{0.2cm}
}


\begin{document}
\maketitle
\subsection*{\noindent Exercise 1.}
QUESTION INFORMATION
\begin{enumerate}[ label=\bfseries(\roman*)]

    \item QUESTION i)
    \begin{proof}\hfill
    \\
	PROOF RIGHT HERE
	\end{proof}
    
    \item QUESTION ii)
    \begin{proof} \hfill
	\\
	PROOF RIGHT HERE
    \end{proof}
\end{enumerate}


\subsection*{\noindent Exercise 2}
QUESTION INFORMATION
\begin{enumerate}[ label=\bfseries(\roman*)]
    
	\item FIRST SUBQUESTION
    \begin{proof} \hfill
	\\
	PROOF RIGHT HERE
    \end{proof}
    
    \item SECOND SUBQUESTION
    \begin{proof}\hfill
	\\
	PROOF RIGHT HERE
    \end{proof}
\end{enumerate}
\end{document}

\end{document}